\documentclass[11pt,a4paper]{article}
\usepackage[margin=22mm]{geometry}
\usepackage{amsmath,amssymb,bm}
\usepackage{booktabs,longtable,array}
\usepackage{enumitem}
\usepackage[hidelinks]{hyperref}
\usepackage{xcolor}
\usepackage[T1]{fontenc}
\usepackage[utf8]{inputenc}

% -----------------------------------------------------------------------------
% SST notation macros. Define notation here; only enter numerical constants when
% the blind protocol explicitly permits them.
% -----------------------------------------------------------------------------
\newcommand{\vswirl}{\mathbf{v}_{\!\boldsymbol{\circlearrowleft}}}
\newcommand{\rhoF}{\rho_{\!f}}
\newcommand{\rhoE}{\rho_{\!E}}
\newcommand{\rhoM}{\rho_{\!m}}
\newcommand{\rc}{r_c}
\newcommand{\SSTTODO}[1]{\ifmmode\text{\color{red}\bfseries TODO}\else\textcolor{red}{\textbf{[REPLACE BEFORE FREEZE: #1]}}\fi}

\title{@@CATALOG-ID@@ --- @@FALSIFIER-NAME@@\\Falsification Protocol and Final Report}
\author{SST-Workbench}
\date{Version @@VERSION@@}

\begin{document}
\maketitle

% SST-REPORT-SECTION:IDENTIFICATION
\section{Identification, scope, and status}
\begin{longtable}{p{0.30\linewidth}p{0.62\linewidth}}
\toprule
Field & Value \\
\midrule
Catalog ID & @@CATALOG-ID@@ \\
Falsifier & @@FALSIFIER-NAME@@ \\
Version & @@VERSION@@ \\
Framework & SST Falsifier Framework v1.0.1 / SST-FALSIFIER-RUNTIME-2 \\
Profile & @@PROFILE@@ \\
Scientific status & \SSTTODO{DRAFT / FROZEN / BLIND-RUN / REVEALED and the applicable date.} \\
\bottomrule
\end{longtable}

\paragraph{Purpose of this document.}
This document is part of the preregistration bundle. It must describe what the code is intended to test \emph{before} target-bearing results are inspected. The frozen SHA-256 bundle therefore covers this file together with the machine-readable science, source, gate, and blindness contracts. After the run, automatically generated provenance and gate results are included without modifying the frozen scientific text.

% SST-REPORT-SECTION:QUESTION
\section{Research question and objective}
\subsection{Single falsifiable research question}
\SSTTODO{Write one precise question. State the independent variable(s), dependent observable(s), admissible system/domain, and the condition that would count against the proposed mechanism.}

\subsection{Scientific objective}
\SSTTODO{Explain what ambiguity this falsifier resolves. Distinguish discovery, confirmation, numerical certification, and physical interpretation. State explicitly what this experiment cannot establish.}

% SST-REPORT-SECTION:HYPOTHESES
\section{Competing hypotheses and allowed conclusions}
\paragraph{Null hypothesis $H_0$.}
\SSTTODO{Give the quantitative null, including expected distribution/range and nuisance parameters.}

\paragraph{Alternative hypothesis $H_1$.}
\SSTTODO{Give the competing quantitative prediction. Avoid encoding private reveal targets into blind discovery.}

\paragraph{Falsification semantics.}
\SSTTODO{State exactly which preregistered gate or statistic falsifies $H_1$, which result fails to reject $H_0$, and which outcomes are merely unresolved. A numerical/backend failure is not physical evidence.}

% SST-REPORT-SECTION:ASSUMPTIONS
\section{Assumptions, domain of validity, and known edge cases}
\begin{enumerate}[label=A\arabic*.]
\item \SSTTODO{Physical assumption; give its domain and how it could fail.}
\item \SSTTODO{Geometric/topological assumption.}
\item \SSTTODO{Numerical/discretization assumption.}
\item \SSTTODO{Statistical independence assumption.}
\end{enumerate}

\paragraph{Known limits.}
\SSTTODO{State small/large parameter limits, continuum/discrete limits, singular limits, and any Newtonian/relativistic or other comparison limits relevant to this falsifier.}

% SST-REPORT-SECTION:SYMBOLS
\section{Symbols, units, and dimensional consistency}
\begin{longtable}{p{0.18\linewidth}p{0.45\linewidth}p{0.14\linewidth}p{0.14\linewidth}}
\toprule
Symbol & Meaning & SI unit & Dimensions \\
\midrule
\SSTTODO{symbol} & \SSTTODO{definition} & \SSTTODO{unit} & \SSTTODO{dimensions} \\
\bottomrule
\end{longtable}

\paragraph{Dimensional audit.}
\SSTTODO{For every nontrivial equation below, show that both sides have the same dimensions. If normalized variables are used, define the dimensional scale that was divided out.}

% SST-REPORT-SECTION:EQUATIONS
\section{Exact mathematical model and observables}
\subsection{Equation register}
Every equation used by the code must have an identifier matching \texttt{science\_contract.json}. Do not write only prose such as ``compute the spectrum'': specify the actual operator, normalization, estimator, or discretization.

\paragraph{E1 --- primary model equation.}
\begin{equation}
\SSTTODO{Insert the exact continuous equation, with all coefficients and boundary/closure assumptions defined.}
\label{eq:E1}
\end{equation}
\noindent Inputs: \SSTTODO{variables/data}. Output: \SSTTODO{observable/intermediate quantity}.\\
Dimensional check: \SSTTODO{show dimensions explicitly}.

\paragraph{E2 --- discrete/numerical form.}
\begin{equation}
\SSTTODO{Insert the actual discrete estimator/operator evaluated by Python/C++/SYCL.}
\label{eq:E2}
\end{equation}
\noindent State the grid/segment convention, indexing, regularization, normalization, quadrature rule, and floating-point precision.

\paragraph{E3 --- decision statistic.}
\begin{equation}
T(D)=\SSTTODO{Define the target-free test statistic as a function of data $D$.}
\label{eq:E3}
\end{equation}

\paragraph{E4 --- null/uncertainty model.}
\begin{equation}
\SSTTODO{Define the null distribution, resampling estimator, likelihood, confidence interval, or error model.}
\label{eq:E4}
\end{equation}

\subsection{Machine-readable equation/step register}
\IfFileExists{AUTO_SCIENCE_CONTRACT.tex}{\input{AUTO_SCIENCE_CONTRACT.tex}}{\emph{AUTO\_SCIENCE\_CONTRACT.tex is generated at FREEZE/run time from science\_contract.json.}}

% SST-REPORT-SECTION:ALGORITHM
\section{Exact algorithm: equation-to-code mapping}
Write the full causal sequence. Every step must say what enters, what mathematical transformation occurs, what implementation performs it, what leaves the step, and which gate consumes it.

\begin{enumerate}[label=S\arabic*.]
\item \SSTTODO{Load source(s); list parser and canonical units. No scientific value is inspected before provenance/admissibility checks.}
\item \SSTTODO{Geometry/topology preprocessing. Give exact formulas/transforms and tolerances.}
\item \SSTTODO{Compute the registered reference observable in Python FP64; cite E-equation IDs.}
\item \SSTTODO{Apply the target-free discovery statistic and null model; state random seed/stopping rule.}
\item \SSTTODO{Run held-out/independent confirmation without retuning discovery choices.}
\item \SSTTODO{Run strict C++ FP64 certification. State what happens if the C++ build is unavailable: certification must fail/unresolve, never silently pass via Python.}
\item \SSTTODO{If enabled, run SYCL screening and CPU--GPU parity. FP32 screening cannot become confirmatory merely because it agrees.}
\item \SSTTODO{Run convergence/refinement and cross-source replication.}
\item \SSTTODO{Only after blind gates, verify reveal commitment and perform target comparison.}
\end{enumerate}

\paragraph{Pseudocode contract.}
\begin{verbatim}
verify_frozen_protocol()
resolve_and_hash_sources()
G0_protocol_integrity()
G1_source_admissibility()
G2_reference_numerics()
G3_target_free_discovery()
G4_independent_confirmation()
G5_strict_cpu_native_certification()
G6_optional_gpu_parity_screening()
G7_cross_source_replication()
G8_mechanism_provider()
G9_reveal_only_after_blind_chain()
\end{verbatim}
\SSTTODO{Modify this pseudocode only if the gate plan is changed before freezing. Explain all deviations.}

% SST-REPORT-SECTION:SOURCES
\section{Source, provenance, and independence contract}
\SSTTODO{List every admissible library/dataset/provider, version/catalog ID, source role (discovery/confirmation), parser/translator, and provenance family. Explain why claimed independent units are genuinely independent and why mesh refinements, seeds, encodings, or reruns are not counted as new samples.}

\paragraph{Exclusion rules.}
\SSTTODO{State objective source exclusion rules before outcome scoring.}

% SST-REPORT-SECTION:GATES
\section{Gate DAG and decision semantics}
The default causal order is
\begin{equation}
G_0\rightarrow G_1\rightarrow G_2\rightarrow G_3\rightarrow G_4\rightarrow G_5\rightarrow G_6\rightarrow G_7\rightarrow G_8\rightarrow G_9.
\end{equation}
A gate may not report \texttt{PASS}, \texttt{FAIL}, or \texttt{UNRESOLVED} unless all declared prerequisites are \texttt{PASS}; otherwise it is \texttt{NOT\_RUN\_PREREQUISITE}. Profile-disabled gates are \texttt{NOT\_APPLICABLE}.

\begin{longtable}{p{0.08\linewidth}p{0.38\linewidth}p{0.22\linewidth}p{0.22\linewidth}}
\toprule
Gate & Scientific question & PASS criterion & FAIL/UNRESOLVED criterion \\
\midrule
$G_0$ & Protocol/runtime integrity & \SSTTODO{criterion} & \SSTTODO{criterion} \\
$G_1$ & Source provenance/admissibility & \SSTTODO{criterion} & \SSTTODO{criterion} \\
$G_2$ & Reference numerical validity & \SSTTODO{criterion} & \SSTTODO{criterion} \\
$G_3$ & Blind discovery & \SSTTODO{criterion} & \SSTTODO{criterion} \\
$G_4$ & Independent confirmation & \SSTTODO{criterion} & \SSTTODO{criterion} \\
$G_5$ & Native CPU certification & \SSTTODO{criterion} & \SSTTODO{criterion} \\
$G_6$ & GPU parity/screening & \SSTTODO{criterion} & \SSTTODO{criterion} \\
$G_7$ & Cross-source replication & \SSTTODO{criterion} & \SSTTODO{criterion} \\
$G_8$ & Mechanism/provider & \SSTTODO{criterion} & \SSTTODO{criterion} \\
$G_9$ & Reveal integrity/comparison & \SSTTODO{criterion} & \SSTTODO{criterion} \\
\bottomrule
\end{longtable}

% SST-REPORT-SECTION:NUMERICS
\section{Numerical certification, convergence, and error budget}
\subsection{Discretization and refinement}
\SSTTODO{Define the refinement sequence, e.g. $N,2N,4N$, time-step relation, fixed final physical time, interpolation/reparameterization rules, and order-of-convergence estimator.}

If three resolutions are used, a typical observed order estimator is
\begin{equation}
 p_{\mathrm{obs}}=\frac{\ln\!\left(\|Q_N-Q_{2N}\|/\|Q_{2N}-Q_{4N}\|\right)}{\ln 2},
\end{equation}
but use this equation only when its assumptions apply; otherwise replace it with the registered estimator.

\subsection{Error budget}
\begin{longtable}{p{0.30\linewidth}p{0.25\linewidth}p{0.35\linewidth}}
\toprule
Error source & Bound/estimator & Acceptance threshold \\
\midrule
Discretization & \SSTTODO{estimator} & \SSTTODO{threshold} \\
Floating-point/backend & \SSTTODO{estimator} & \SSTTODO{threshold} \\
Source/parser & \SSTTODO{estimator} & \SSTTODO{threshold} \\
Statistical & \SSTTODO{estimator} & \SSTTODO{threshold} \\
\bottomrule
\end{longtable}

% SST-REPORT-SECTION:BACKENDS
\section{Backend authority and parity}
The backend contract must distinguish \emph{requested} from \emph{actual} execution. For candidate vector $\mathbf q_b$ and authoritative reference $\mathbf q_r$, the default parity diagnostic is
\begin{equation}
\epsilon_{b/r}=\frac{\|\mathbf q_b-\mathbf q_r\|_2}{\max(\|\mathbf q_r\|_2,\epsilon_0)}.
\end{equation}
\SSTTODO{State the registered $\epsilon_0$, C++/Python tolerance, GPU/CPU tolerance, dtype of each backend, and whether each backend has REFERENCE, CERTIFICATION, SCREENING\_ONLY, or DEVELOPMENT authority.}

\paragraph{No silent fallback.}
If a gate requires C++/OpenMP or SYCL and the actual backend is Python, that backend gate cannot pass. A fallback may keep development code runnable but has no certification authority.

% SST-REPORT-SECTION:BLINDNESS
\section{Blindness, opaque identifiers, and reveal}
Opaque identifiers use HMAC-SHA256 with a private key, not a public salt:
\begin{equation}
\mathrm{ID}=\operatorname{Trunc}_n\!\left(\operatorname{HMAC}_{\mathrm{SHA256}}(K_{\mathrm{private}},\mathrm{label})\right).
\end{equation}
The reveal commitment is nonced:
\begin{equation}
C=\operatorname{SHA256}(\texttt{SST-REVEAL-COMMIT-2}\,\|\,N_{\mathrm{private}}\,\|\,R),
\end{equation}
where $R$ is the exact reveal payload and $N_{\mathrm{private}}$ is disclosed only when reveal is allowed.

\SSTTODO{List what information is blinded, who/what can access the private mapping, which terms are forbidden in blind artifacts, and the exact condition permitting reveal.}

% SST-REPORT-SECTION:STATISTICS
\section{Statistical/null model and multiplicity}
\SSTTODO{Define the null ensemble, randomization/bootstrap/permutation procedure, number of trials or stopping rule, confidence level, effect-size measure, multiple-testing correction, and the independent sampling unit.}

\paragraph{Selection effects.}
\SSTTODO{Explain how hyperparameter search, knot/library screening, seed selection, and post-hoc thresholds are prevented from leaking into confirmation.}

% SST-REPORT-SECTION:RESULTS
\section{Results and automatic run provenance}
This section is populated after execution. The frozen scientific text above must not be rewritten to match the outcome.
\IfFileExists{AUTO_RESULTS.tex}{\input{AUTO_RESULTS.tex}}{\emph{AUTO\_RESULTS.tex is generated after a run.}}

\subsection{Primary numerical result}
The numerical result, uncertainty, gate status, protocol hash, backend identity, and source provenance are inserted by \texttt{AUTO\_RESULTS.tex}. The frozen protocol text is never rewritten to match the outcome.

% SST-REPORT-SECTION:INTERPRETATION
\section{Interpretation, limits, falsifiable predictions, and minimal follow-up experiments}
\paragraph{Preregistered interpretation rule.}
\SSTTODO{State in advance how PASS, FAIL, UNRESOLVED, backend failure, and source failure will be interpreted without moving thresholds after seeing the result.}

\paragraph{Preregistered uncertainty/edge-case rule.}
\SSTTODO{List the dominant anticipated uncertainty sources and conditions under which the planned conclusion would not be valid.}

\paragraph{Falsifiable predictions reserved for independent checks.}
\SSTTODO{List predictions not used for fitting/selection and the experiment that could falsify each prediction.}

\subsection{Post-run discussion}
\IfFileExists{POST_RUN_DISCUSSION.tex}{\input{POST_RUN_DISCUSSION.tex}}{\emph{Optional post-run discussion is kept outside the frozen protocol in POST\_RUN\_DISCUSSION.tex.}}

% SST-REPORT-SECTION:REPRODUCIBILITY
\section{Reproducibility and exact commands}
\begin{verbatim}
run_all.cmd FREEZE
run_all.cmd BASIC
run_all.cmd FULL
run_all.cmd CERTIFY
run_all.cmd REVEAL
\end{verbatim}
\SSTTODO{Add any required Workbench root, source overrides, random seeds, compiler/oneAPI setup, thread counts, and hardware constraints.}

\section{References}
For every external/non-original equation, comparison model, numerical method, or dataset definition used by the falsifier, add a complete \verb|\bibitem| entry with authors, year, title, venue, and DOI/permalink where available.

\begin{thebibliography}{99}
\bibitem{replace_reference}
\SSTTODO{Authors (year), ``Title,'' Venue, DOI/permalink.}
\end{thebibliography}

\end{document}
